\section{Euclidean space, finite volume, and the five solid types}
\label{sec:solids}
In $\K^d$ let $u\cdot v=\sum_{i<d}u_iv_i$ and
$N(u)=\sum_{i<d}u_i^2$. Dimension $d$ may be a natural-number input
for these bounded sums. In dimension three write
\[
 u\times v=(u_2v_3-u_3v_2,\ u_3v_1-u_1v_3,\ u_1v_2-u_2v_1),
 \qquad [u,v,w]=u\cdot(v\times w).
\]
All surfaces and solids below have finite vertex, face, or triangulation
data. The section addresses the main article's solid-geometry, scaling,
and Platonic-solid statements \cite{wiki,euclid,platoniccore}.

\begin{theorem}[Lagrange's identity in finite dimension and cross products]
\label{sol:gram}
For every finite $d$,
\[
 N(u)N(v)-(u\cdot v)^2=
       \sum_{i<j<d}(u_iv_j-u_jv_i)^2\ge0.
\]
Consequently Cauchy--Schwarz and the Euclidean metric triangle inequality
hold in every such dimension. In three dimensions the sum on the right
is $N(u\times v)$, and $u\times v$ is perpendicular to both factors.
\end{theorem}
\begin{proof}
Expand the two products on the left. The diagonal terms cancel. Group
the remaining terms with indices $(i,j)$ and $(j,i)$; their sum is
$u_i^2v_j^2+u_j^2v_i^2-2u_iv_iu_jv_j$, the asserted square.
This reindexing is a bounded-sum identity, proved by induction on $d$
with all coordinate lists as parameters. The inequality follows because
squares are nonnegative. The square-root and norm-expansion argument in
Theorem~\ref{tri:metric} now proves Cauchy--Schwarz and the triangle
inequality word for word. In dimension three the three squared terms
are the squared coordinates of the cross product. Its dot products with
$u,v$ vanish by cancellation of their six terms.
\end{proof}

\begin{theorem}[Planes and exact intersections]\label{sol:planes}
Three noncollinear spatial points determine a unique plane. A plane
$n\cdot X=h$, $N(n)>0$, and line $A+tu$ meet uniquely if
$n\cdot u\ne0$, at
\[
 A+\frac{h-n\cdot A}{n\cdot u}u.
\]
Otherwise the whole line lies in the plane when $n\cdot A=h$, and
the line is disjoint from it when that equality fails. Two nonparallel
planes intersect in a line with explicitly computable coefficients.
\end{theorem}
\begin{proof}
For points $A,B,C$ take $n=(B-A)\times(C-A)\ne0$ and $h=n\cdot A$.
The cross-product calculation gives all three incidences. The normal
space perpendicular to the two independent vectors has dimension one:
solve the two equations using a nonzero two-by-two minor, which exists
because the cross product is nonzero. Thus any other plane through the
three points has proportional normal and the same equation up to scale.
The line cases follow by solving $n\cdot A+t(n\cdot u)=h$.

For planes $n\cdot X=h$, $m\cdot X=k$, assume
$D=N(n)N(m)-(n\cdot m)^2=N(n\times m)>0$. Put
\[
 X_0=\frac{hN(m)-k(n\cdot m)}D n
       +\frac{kN(n)-h(n\cdot m)}D m.
\]
Taking its two dot products yields $h,k$. Every point
$X_0+t(n\times m)$ satisfies the equations. Subtracting $X_0$ from
any other solution gives a vector perpendicular to $n,m$, hence a
multiple of their cross product by the same nonzero-minor argument.
The parallel-normal case is decided by the compatibility of the constants.
\end{proof}

\begin{theorem}[Distance to a plane and orthogonal projection]
\label{sol:projection}
The foot of $P$ on $n\cdot X=h$ is
$F=P+(h-n\cdot P)n/N(n)$, and its distance is
$|n\cdot P-h|/\sqrt{N(n)}$. It uniquely minimizes the distance to
$P$ among points of the plane. Reflection $P\mapsto2F-P$ is an
involutive isometry fixing exactly the plane.
\end{theorem}
\begin{proof}
Substitution verifies $n\cdot F=h$. The vector $P-F$ is a scalar
multiple of $n$, while every displacement $X-F$ in the plane is
perpendicular to $n$. Thus
$N(P-X)=N(P-F)+N(X-F)$, proving the distance and uniqueness claims.
Reflection negates the normal component while preserving the component
in the plane. The orthogonal sum of their squared norms proves isometry,
and the sign reversal proves involution and the fixed-point assertion.
\end{proof}

\begin{theorem}[Sphere sections and tangency]\label{sol:sphere-section}
Let a sphere have center $O$ and radius $R>0$, and let a plane be at
distance $d$ from $O$, with foot $F$. Its intersection with the sphere
is empty if $d>R$, the singleton $F$ if $d=R$, and the circle in
that plane centered at $F$ with squared radius $R^2-d^2$ if $d<R$.
\end{theorem}
\begin{proof}
For $X$ in the plane, orthogonal decomposition gives
$N(X-O)=N(X-F)+d^2$. The sphere equation is therefore
$N(X-F)=R^2-d^2$. A negative right side has no solution; zero
forces $X=F$. In the positive case choose an orthonormal frame in
the plane by normalizing any nonzero perpendicular vector and its
cross product with the normal. The equation is exactly the planar
circle equation. A suitable perpendicular vector is obtained from a
coordinate axis not parallel to $n$, a three-case decidable selection.
\end{proof}

\begin{theorem}[A sphere through four noncoplanar points]\label{sol:circumsphere}
Four noncoplanar points determine a unique sphere. With one point at
zero and the others $u,v,w$, let $D=[u,v,w]\ne0$. Its center is
\[
 O=\frac{N(u)(v\times w)+N(v)(w\times u)+N(w)(u\times v)}{2D}.
\]
Its squared radius is $N(O)>0$.
\end{theorem}
\begin{proof}
Taking dot products gives $2O\cdot u=N(u)$ and the corresponding
equations for $v,w$. Expanding $N(O-u)$, $N(O-v)$, and $N(O-w)$
therefore gives $N(O)$. Conversely any sphere through zero and those
points must satisfy these three linear equations. Their determinant
is nonzero, so the solution is unique. If $O=0$, the first equation
would force $u=0$, contradicting noncoplanarity, hence the radius is positive.
\end{proof}

\begin{theorem}[Finite oriented volume and its boundary formula]
\label{sol:volume}
For a tetrahedron $ABCD$ define its signed volume by
\[
 V_s(A,B,C,D)=\frac{[B-A,C-A,D-A]}6.
\]
For a finite positively oriented conforming tetrahedral partition, the sum of these
volumes depends only on its oriented triangular boundary, and is additive
under finite dissection. An affine map with linear part $T$ multiplies
signed volume by $\det T$ and unsigned volume by $|\det T|$.
In particular a similarity of ratio $\lambda$ multiplies volume by
$\lambda^3$.
Here conforming means that shared boundary pieces have identical face
triangulations; otherwise first take their finite common refinement.
\end{theorem}
\begin{proof}
Multilinearity gives the four-face identity
\[
 [B-A,C-A,D-A]=[B,C,D]-[A,C,D]+[A,B,D]-[A,B,C].
\]
Sum this over the tetrahedra. Every internal oriented triangular face
occurs twice with opposite signs; only the boundary remains. This proves
independence, additivity, and invariance under changing an interior
triangulation. The sum can be updated one face at a time with a
quantifier-free equality invariant.

The permutation formula for the determinant gives
$[Tu,Tv,Tw]=\det(T)[u,v,w]$. A common translation disappears from
the edge differences, proving the affine claim. If $T=\lambda Q$ and
$Q^\mathsf TQ=I$, then $(\det Q)^2=1$, so $|\det T|=\lambda^3$.
For a unit box the familiar six tetrahedra obtained by ordering its three
coordinates each have determinant one; their volumes sum to one.
Thus this volume has the unit-box normalization as well as finite
dissection additivity. It is the ordinary polyhedral volume on these
finite representations, not a measure on arbitrary subsets.
\end{proof}

\begin{theorem}[Prism and box volumes]\label{sol:prism}
A prism on a planar triangulated base of area $A$ and perpendicular
height $h$ has volume $Ah$. In particular a rectangular box has the
product of its three edge lengths as volume.
\end{theorem}
\begin{proof}
First take a triangular base $A_0B_0C_0$ and displacement $u$.
Its prism divides into tetrahedra with vertex sets
\[
 \{A_0,B_0,C_0,C_0+u\},\quad
 \{A_0,B_0,B_0+u,C_0+u\},\quad
 \{A_0,A_0+u,B_0+u,C_0+u\}.
\]
In base barycentric coordinates and height parameter $0\le t\le1$,
these three tetrahedra are the three consecutive regions obtained by
ordering $t$ against the two cumulative barycentric coordinates. They
cover the prism and have disjoint interiors. Each has absolute determinant
$|[B_0-A_0,C_0-A_0,u]|$. Their total volume is half that determinant.
The base area is $\sqrt{N((B_0-A_0)\times(C_0-A_0))}/2$ and
the height is the absolute normal component of $u$, so their product
is the same half-determinant. For a base with several triangles choose
one global vertex order before applying these splits, so that shared
vertical rectangles acquire the same diagonal. Sum over the base triangulation; internal
vertical faces cancel. A rectangular base has the product of its two
perpendicular edges as area, giving the box formula.
\end{proof}

\begin{theorem}[Pyramid volume and the one-third ratio]\label{sol:pyramid}
A pyramid on a planar triangulated base of area $A$, with apex at
perpendicular height $h$, has volume $Ah/3$. It therefore has one third
the volume of a prism with the same base and height.
\end{theorem}
\begin{proof}
Each base triangle together with the apex is a tetrahedron. Its absolute
triple product is its doubled base area times the perpendicular height,
because only the normal component of the apex displacement contributes.
Divide by six to obtain one third the base area times $h$.
All base triangles have the same height, and their areas add, giving
$Ah/3$. Their interiors and conical sectors are disjoint, so the finite
volume additivity applies. The prism comparison follows from
Theorem~\ref{sol:prism}. The height-zero case gives zero on both sides.
\end{proof}

\begin{theorem}[Euler's formula for a finite convex polyhedron]
\label{sol:euler}
For a three-dimensional convex polyhedron given by its exact finite
vertices and maximal polygonal faces, with counts $V,E,F$,
$V-E+F=2$.
\end{theorem}
\begin{proof}
Choose a face $F_0$ and its vertex centroid $p$, which is in its relative
interior. Let its supporting equation be $n\cdot X=h$, with the solid
on the $\le h$ side. Choose a viewpoint
$z=p+\epsilon n/N(n)$, where $\epsilon>0$ is sufficiently small.
This choice is finite and effective. At $p$ all other face inequalities
are strict, so their finitely many positive gaps give a bound ensuring
that $z$ remains inside those halfspaces. For every vertex $v$ not in
$F_0$, let $\delta_v=h-n\cdot v>0$. Its projection from $z$ to
the face plane is
\[
 p+\frac{\epsilon}{\delta_v+\epsilon}(v_{\parallel}-p),
 \qquad
 v_{\parallel}=v+\frac{h-n\cdot v}{N(n)}n.
\]
Each of the finitely many edge inequalities of $F_0$ is strict at $p$.
Choose $\epsilon$ below the minimum of the resulting positive rational
bounds after bounding the displayed displacements. Then every such
projection is strictly inside $F_0$.
More explicitly, let $g_i>0$ be rational lower bounds for the gaps to
the other face planes, and choose rational
$A_i\ge|n_i\cdot n/N(n)|$. For each edge inequality
$m_j\cdot(X-p)+b_j\ge0$ of $F_0$, choose a rational
$g_j>0$ below $b_j$, and choose rational
$d_v>0$ below $\delta_v$ and
$L_{jv}\ge|m_j\cdot(v_{\parallel}-p)|$.
It suffices to take
\[
 \epsilon=\min\left(1,\min_i\frac{g_i}{2(1+A_i)},
                     \min_{j,v}\frac{d_vg_j}{2(1+L_{jv})}\right).
\]
All minima are finite and positive; omit an empty subfamily. The two
changes in the relevant inequalities are then smaller than half their
strict margins.

Every ray from $z$ through the face meets the convex solid in an interval:
each face inequality cuts out an interval of ray parameters, and their
finite intersection is an interval. Its far endpoint is unique. Hence the
remaining surface projects onto $F_0$ without crossing or overlap of
face interiors. Projection denominators are positive, so a convex
combination projects to a convex combination after positive renormalization
of its weights. Each projected face is therefore convex.
Triangulate each projected convex face by a fan. We now
have a finite triangulated disk with convex polygonal boundary. Let $b$
be its number of boundary vertices, $V_i$ its number of interior vertices,
$E_i$ its number of interior edges, and $T$ its number of triangles.

At each interior vertex the incident triangle sectors fill a full turn;
at boundary vertices they fill the outer polygon's interior angles.
The triangle angle-sum theorem and Theorem~\ref{poly:angles}, using
their finite half-turn carries, therefore give
$T=2V_i+b-2$. Counting triangle-edge incidences gives $3T=2E_i+b$.
Thus, with $V=V_i+b$ and $E=E_i+b$,
$V-E+T=1$. Add back the omitted outer face to get two. Every fan
diagonal that was introduced adds one edge and one face, so removing
them preserves $V-E+F$. This proves the original formula.

All choices are from finite vertex and face lists, and every strict
geometric margin has an algebraic sign gap. The counts and incidence
equalities are bounded folds. The argument uses neither a theorem about
arbitrary topological surfaces nor an unbounded compactness principle.
\end{proof}

\begin{theorem}[Only five regular face--vertex types]
\label{sol:five-types}
Suppose a convex polyhedron has regular $p$-gonal faces with $q$ faces
meeting at every vertex, where $p,q\ge3$. Then
\[
 (p,q)\in\{(3,3),(4,3),(3,4),(5,3),(3,5)\}.
\]
Its counts $(V,E,F)$ are respectively
$(4,6,4),(8,12,6),(6,12,8),(20,30,12),(12,30,20)$.
\end{theorem}
\begin{proof}
Face--edge counting gives $pF=2E$, while face--vertex counting gives
$qV=pF=2E$. Multiply Euler's equality
by $pq$ and substitute to obtain
\[
 E(2p+2q-pq)=2pq>0.
\]
Hence $(p-2)(q-2)<4$. Both factors are positive integers. Enumerating
their products below four gives $(1,1),(1,2),(2,1),(1,3),(3,1)$,
which are exactly the listed possibilities. Finally
$E=2pq/(2p+2q-pq)$ and the incidence equations give the displayed
counts. This is also the arithmetic core already represented by
\texttt{Geometry/PlatonicSolids.lean} in ProveIt \cite{platoniccore}.
\end{proof}

\begin{theorem}[Exact realizations of all five types]\label{sol:realizations}
All five types in Theorem~\ref{sol:five-types} are realized by explicit
convex polyhedra with constructible coordinates. Put
$\varphi=(1+\sqrt5)/2$, so $\varphi^2=\varphi+1$.
Their vertices may be chosen as follows:
\begin{center}\small
\begin{tabular}{ll}
\toprule
Solid & Vertex set\\
\midrule
Tetrahedron & $(\pm1,\pm1,\pm1)$ with sign product $+1$\\
Cube & all $(\pm1,\pm1,\pm1)$\\
Octahedron & $(\pm1,0,0),(0,\pm1,0),(0,0,\pm1)$\\
Icosahedron & cyclic permutations of $(0,\pm1,\pm\varphi)$\\
Dodecahedron & cube vertices and cyclic permutations of
$(0,\pm\varphi^{-1},\pm\varphi)$\\
\bottomrule
\end{tabular}
\end{center}
\end{theorem}
\begin{proof}
For the first three sets, the supporting planes and faces follow directly
from their sign patterns. The tetrahedron has four equilateral faces
of squared edge length 8; the cube has six squares of squared edge
length 4; the octahedron has eight equilateral triangles of squared
edge length 2. Their vertex squared radii are $3,3,1$, respectively.

For the icosahedron take outward normals of two kinds: all eight signed
$(1,1,1)$, and the twelve cyclic permutations of
$(0,\pm\varphi,\pm\varphi^{-1})$. Each supporting value is
$\varphi^2$. Substitution of the twelve vertices gives at most this
value, with exactly three equalities. Representatives of the equality
triples are
\[
 (0,1,\varphi),(1,\varphi,0),(\varphi,0,1)
\quad\hbox{and}\quad
 (0,1,\varphi),(1,\varphi,0),(-1,\varphi,0).
\]
All three pairwise distances in either triple are 2. Signed coordinates
and cyclic permutations give the twenty facets. Every vertex has five
incident triangles and all vertex squared radii are $\varphi+2$.

For the dodecahedron use the twelve cyclic signed normals
$(0,\pm\varphi,\pm1)$ with supporting value $\varphi^2$.
For $n=(0,\varphi,1)$ the five equality vertices, in boundary order, are
\[
 (1,1,1),\quad (\varphi^{-1},\varphi,0),\quad
 (-\varphi^{-1},\varphi,0),\quad (-1,1,1),\quad
 (0,\varphi^{-1},\varphi).
\]
They have common squared radius 3 and lie in one plane, so they lie
on the circle with center the foot of the origin in that plane.
Adjacent squared distances are $4/\varphi^2=6-2\sqrt5$, and
successive turns have the same positive sign and the same dot product.
Thus the face is a convex regular pentagon. The sign and cyclic
symmetries give twelve such faces, with three meeting at each vertex.

For completeness the accompanying exact certificate enumerates every
noncollinear vertex triple, retains the plane precisely when all remaining
vertices are on one side, and merges equal equality sets. This exhausts
the facets: a two-dimensional exposed face of a finite convex hull is
the convex hull of its equality vertices and contains three affinely
independent ones. All comparisons take place in $\Q(\sqrt5)$ with
integer sign tests. The returned facets are exactly the ones above;
the full oriented face and edge lists are in
\texttt{certificates/solids.json}. Their finite counts and face equations
are reproducible by \texttt{code/check\_solids.py}.
\end{proof}

\paragraph{Exact extent of the classification.}
The last two theorems prove the five admissible face--vertex types and
give regular convex examples of each, including their incidence counts.
They do not add a separate global rigidity or abstract-combinatorial
uniqueness theorem asserting that every polyhedron with a given type is
congruent to the displayed one. The package's inventory records this
distinction explicitly. The exact coordinate certificates establish
existence and regularity, rather than relying on a drawing of a solid.
